\documentclass[11pt]{article}
\usepackage[margin=1in]{geometry}
\usepackage{amsmath,amssymb,amsthm,booktabs}
\newtheorem{theorem}{Theorem}
\newtheorem{proposition}{Proposition}
\newtheorem{example}{Example}
\title{State-Weighted Defect Transfer for Dependent Block Explorations}
\author{}
\date{}
\begin{document}
\maketitle

\begin{abstract}
We prove a weighted positive-kernel criterion for exponential suppression of
connected bad-block events.  The criterion remembers energetic cost stored in
a retained boundary state and can remain subcritical even when a scalar
worst-history conditional probability equals one.  It is an abstract finite
transfer theorem; construction of the required Yang--Mills kernel remains an
analytic obligation.
\end{abstract}

\section{Weighted exploration}
Let $\Sigma$ be a measurable retained-state space, $V:\Sigma\to[1,\infty)$,
and let $Q_i(s,ds')$ be the subprobability kernel for a bad extension at step
$i$.  Assume uniformly in $i$ and $s$ that
\begin{equation}\label{eq:drift}
 \int V(s')Q_i(s,ds')\le qV(s).
\end{equation}

\begin{theorem}[Weighted chain bound]
If \eqref{eq:drift} holds and the initial law $\nu$ satisfies
$\nu(V)\le M$, then any specified chain of $n$ bad extensions has probability
at most $Mq^n$.
\end{theorem}
\begin{proof}
Positivity and repeated kernel composition give
\[
 \int \nu(ds_0)Q_1(s_0,ds_1)\cdots Q_n(s_{n-1},ds_n)V(s_n)
 \le q^n\nu(V)\le Mq^n.
\]
Since $V(s_n)\ge1$, deleting the terminal weight can only decrease the
integrand on the bad-chain event.
\end{proof}

\begin{theorem}[Rooted connected-set sum]
On a graph of maximum degree $\Delta$, suppose every connected bad set can be
exposed by successive adjacent bad extensions satisfying the preceding
hypothesis.  If $e^\alpha\Delta^2q<1$, then
\[
 \sum_{\substack{X\ni o\\X\text{ connected}}}
 e^{\alpha|X|}\mathbb P(B_X)
 \le \frac{M e^\alpha q}{1-e^\alpha\Delta^2q}.
\]
\end{theorem}
\begin{proof}
A rooted connected set of size $n$ admits a depth-first traversal of length at
most $2(n-1)$, giving the standard upper bound
$\Delta^{2(n-1)}$ on its number.  Sum the weighted chain estimate over $n$.
\end{proof}

\section{Why the state weight matters}
\begin{example}
For $0<q<1$, take
\[
 Q=\begin{pmatrix}0&q^2\\1&0\end{pmatrix},\qquad
 V=\begin{pmatrix}1\\q^{-1}\end{pmatrix}.
\]
The scalar row ceiling is one, but $QV=qV$.  Entry into the high-cost state is
rare, although the following bad transition is forced.  For finite irreducible
$Q$, minimizing $\max_s(QV)(s)/V(s)$ over positive $V$ yields the
Perron--Frobenius spectral radius.
\end{example}

At $q=0.002430277501841998$, $\Delta=18$, and $\alpha=0.1$, the branching
ratio is $0.8702225338$ and the rooted sum is $0.0206959813$.  The strict
critical value is
\[
 q_{\rm crit}=e^{-0.1}/18^2=0.0027927080803578997,
\]
equivalently a weighted extension cost greater than $5.8807435158$.

\section{Scope}
The theorem does not identify $Q$ with a nonlinear lattice Yang--Mills
conditional kernel.  The next obligation is to construct a gauge-covariant
retained state and physical weight satisfying the strict ceiling uniformly in
volume.  No RG, continuum, or Hamiltonian-gap claim follows at this stage.
\end{document}

